\section{Method: From Question to Full-Marks Answer}\label{sec:method}

\subsection{The 8 steps (inside each part of the question)}
{\small\renewcommand{\arraystretch}{1.3}
\begin{longtable}{L{0.6cm}L{2.6cm}L{9.0cm}L{3.4cm}}
\toprule
\tblhead{\textbf{\#} & \textbf{Step} & \textbf{What to do} & \textbf{Where}}\\
\midrule
\endhead
1 & \textbf{Givens} & List architecture (cores/PEs, issue width, in-order/OoO, threads, caches, memory, interconnect), workload (operations, sizes, access pattern, dependences, branches), symptom (with its number). & Breakdown grid \pg{sec:breakdown} \\
2 & \textbf{Assumptions} & Numbered A1, A2\ldots{} for anything missing; few and plausible. & Typical values \pg{sec:typical} \\
3 & \textbf{Workload analysis} & Which parallelism (ILP/DLP/TLP) and how much; \emph{order of execution} (dependences, critical path, serial part, reductions); memory behaviour (working set vs.\ cache, reuse, stride); sharing/communication; branches. Write ``parallelism present: yes/no, where''. & Order of execution \pg{grp:work}; memory test \pg{sec:memworks} \\
4 & \textbf{Diagnosis} & Root cause from the Symptom Index. Prove with one number; rule out one rival explanation. & Symptom Index \pg{sec:symidx} \\
5 & \textbf{Redesign every layer} & One concept per bottleneck, for every layer the question touches. & Layers \pg{sec:layers}; Clue Index \pg{sec:clueidx} \\
6 & \textbf{Static or dynamic} & Tag each choice and say why it suits this workload. & \pg{sec:staticdyn} \\
7 & \textbf{Justify} & ``Why it beats'' column vs.\ $\ge$2 alternatives; weakness accepted; ``Needs'' added. & Concept tables; knock-on \pg{sec:knockon} \\
8 & \textbf{Outcome} & Expected gain and where the bottleneck moves next. & Model answers \pg{sec:modelA} \\
\bottomrule
\end{longtable}}

\subsection{Answer layout (copy these headings)}\label{sec:layout}
{\small\renewcommand{\arraystretch}{1.3}
\begin{longtable}{L{4.2cm}L{9.0cm}L{2.4cm}}
\toprule
\tblhead{\textbf{Heading} & \textbf{What goes under it} & \textbf{Time share}}\\
\midrule
\endhead
1. Given \& assumptions & Three lists (architecture, workload, symptom) + numbered assumptions & 10\% \\
2. Workload analysis & Parallelism kind and amount, order of execution, memory behaviour, communication & 15\% \\
3. Why it is not ideal & Mechanism + one number + one rival explanation ruled out & (with 2) \\
4. Proposed architecture & Table: \emph{Layer $|$ Change $|$ S/D $|$ Why for this workload}; one small diagram & 45\% \\
5. Why this is better & Each choice vs.\ $\ge$2 alternatives, with numbers; one lab cited & 20\% \\
6. Demerits and next bottleneck & Costs accepted; which limit binds after the fix; final checklist \pg{sec:checklist} & 10\% \\
\bottomrule
\end{longtable}}

\subsection{Make every answer specific, not generic}\label{sec:specific}
{\small\renewcommand{\arraystretch}{1.3}
\begin{longtable}{L{0.6cm}L{3.6cm}L{11.4cm}}
\toprule
\tblhead{\textbf{\#} & \textbf{Rule} & \textbf{What it looks like}}\\
\midrule
\endhead
R1 & Quote the given & ``the 64\,MB table is 32$\times$ the 2\,MB L2'' \\
R2 & State the mechanism & What physically happens on each access or instruction (``each hop's address depends on the previous load'') \\
R3 & Quantify once & AMAT, CPI, Amdahl, bandwidth or a timeline \\
R4 & Name the exact variant & ``4K-entry gshare + BTB + return stack'', not ``a predictor'' \\
R5 & Reject two alternatives & With the reason each fails \emph{in this scenario} \\
R6 & Accept a cost & And say why it is acceptable here \\
R7 & Add the knock-on & What else must change (``Needs'') \\
R8 & Static or dynamic & And why for this workload \\
\bottomrule
\end{longtable}}

\begin{ex}[Generic vs.\ specific]
\textbf{Generic (few marks):} ``Use multithreading to hide memory latency.''\\[3pt]
\textbf{Specific (full marks):} ``P2 follows hash chains in a 64\,MB table, 32$\times$ the 2\,MB L2, so nearly every hop is a DRAM miss ($\approx$200 cycles, A3), and each hop's address depends on the previous load, so neither a larger OoO window nor more MSHRs can overlap them (MLP $=1$; Lab 4: more MSHRs gave no gain on a dependent chase). I add 4 hardware threads per core with \emph{switch-on-miss} (coarse-grained) multithreading \cp{T1}: stalls are long, so a few-cycle refill is negligible against 200 cycles, and three other independent lookups run meanwhile ($\approx$4$\times$ busy cycles). I reject SMT because these cores are in-order and SMT's low cost relies on OoO renaming \cp{T3}, and a bigger L2 because 64\,MB will not fit on chip. Cost: 4 PCs and register files, more L1 pressure, so I raise L1 to 4-way \cp{M4}. Dynamic: the hardware switches on actual misses.''
\end{ex}

\subsection{What the professor's pointers mean}
{\small\renewcommand{\arraystretch}{1.3}
\begin{longtable}{L{4.2cm}L{8.4cm}L{3.0cm}}
\toprule
\tblhead{\textbf{He said} & \textbf{You write} & \textbf{Where}}\\
\midrule
\endhead
Mention assumptions & Numbered list, few and plausible & \pg{sec:typical} \\
Why is it occurring & Mechanism + number & \pg{sec:symidx} \\
Static / dynamic decision & Tag each choice, with the reason & \pg{sec:staticdyn} \\
SIMT, DLP & Warps, masks, coalescing, occupancy & D table \pg{grp:dlp} \\
Multiple PEs, distribution, timing & Split, timeline, $T_{par}$ & W table \pg{grp:work} \\
Benefits and demerits & Both, always & \pg{sec:paraproscons} \\
Order of execution & Dependence graph, critical path & \pg{grp:work} \\
Parallelism present or not & Fraction $f$, Amdahl & \pg{grp:work} \\
Data, memory, communication & Roofline, memory test, communication cost & \pg{sec:memworks} \\
Synchronisation & What sync, its cost, how to reduce it & C table \pg{grp:coh} \\
Why your solution is better & ``Why it beats'' column, $\ge$2 rivals & Concept tables \\
Lab aspects & One lab finding as evidence & \pg{sec:labs} \\
Adjust the whole architecture & Every layer + knock-on & \pg{sec:layers}, \pg{sec:knockon} \\
\bottomrule
\end{longtable}}

\subsection{Phrase bank for justification}\label{sec:phrases}
{\small\renewcommand{\arraystretch}{1.3}
\begin{longtable}{L{3.6cm}L{12.0cm}}
\toprule
\tblhead{\textbf{To say\ldots} & \textbf{Write}}\\
\midrule
\endhead
Diagnosis & ``The clue `[\ldots]' shows that [\ldots], so the bottleneck is [\ldots], not [rival], because [number].'' \\
Why better & ``I choose X over Y because, in this workload, Y [fails by mechanism] when [property of this scenario], whereas X [mechanism].'' \\
Static / dynamic & ``This is a \emph{static} decision: [\ldots] is known before run time.'' / ``\emph{Dynamic}: [\ldots] is only known at run time.'' \\
Cost accepted & ``The cost is [\ldots]; it is acceptable here because [constraint].'' \\
Next bottleneck & ``After this change the binding limit becomes [\ldots]; I address it with [\ldots].'' \\
Lab evidence & ``This matches Lab [n], where [\ldots].'' \\
Depth (sensitivity) & ``A[n] assumed [\ldots]; if instead [\ldots], I would choose [\ldots].'' \\
\bottomrule
\end{longtable}}

\subsection{Worked breakdown of a big multi-part question}\label{sec:breakdown}
\textbf{Question (practice).} A 4-core chip (2-issue \emph{in-order} cores, one thread each; private 32\,KB direct-mapped L1; shared 2\,MB L2; snooping bus; one DRAM channel 12.8\,GB/s) runs a video-analytics service. \textbf{P1} a filter applies the same arithmetic to every pixel of each 1080p frame, with an \texttt{if} for edge pixels: it reaches 15\% of FP peak. \textbf{P2} a lookup stage follows linked hash chains in a 64\,MB table: cores are stalled most cycles. \textbf{P3} all cores increment shared per-category counters: 4 cores are slower than 2. \textbf{P4} per-frame tasks take 1--10\,ms and are split equally in advance: some cores sit idle. Redesign the processor.

\small
\begin{longtable}{L{0.6cm}L{3.6cm}L{4.3cm}L{3.4cm}L{2.9cm}}
\toprule
\tblhead{\textbf{Part} & \textbf{Clue words (underline)} & \textbf{Concept / root cause} & \textbf{Open concepts} & \textbf{S/D}}\\
\midrule
\endhead
P1 & ``same arithmetic on every pixel'', ``\texttt{if} for edges'', ``2-issue in-order'', ``15\% of peak'' & DLP not exploited by a scalar 2-issue core; branch per pixel; check bandwidth (roofline) & \cp{D2} SIMD (or \cp{D3} GPU), \cp{D5} masks, \cp{D4} layout, \cp{M11} prefetch & Static vectorisation; dynamic masks \\
P2 & ``linked chains'', ``64\,MB'' vs.\ ``2\,MB L2'', ``stalled'' & Dependent DRAM misses, MLP $=1$, capacity far exceeded & \cp{T1}/\cp{T2} multithreading, \cp{M13} layout, \cp{M4} & Dynamic thread switch \\
P3 & ``all cores increment shared counters'', ``4 slower than 2'' & True + false sharing: coherence ping-pong on counter lines & \cp{W5} privatise + reduce, \cp{C3} pad, \cp{C1} & Static \\
P4 & ``1--10\,ms'', ``split equally in advance'', ``idle'' & Load imbalance from static split of unequal tasks & \cp{W3} dynamic queue, chunk size & Dynamic \\
All & ``one DRAM channel'', ``snooping bus'' & SIMD + 4 threads/core raise memory demand & \cp{M18} banks/channels, \cp{M10} non-blocking & --- \\
\bottomrule
\end{longtable}
\normalsize
\textbf{Full answer: Model G (\pg{sec:modelG}).} \textbf{Then connect them (whole architecture):} SIMD units \emph{and} 4 threads per core both raise memory traffic, so add a second DRAM channel and a non-blocking L1 (knock-on, \pg{sec:knockon}); multithreading on in-order cores means switch-on-miss or fine-grained, not SMT; privatised counters remove the bus traffic that would otherwise grow with the extra threads; the dynamic queue for P4 needs a cheap atomic fetch-and-add on a single queue index \cp{C7}. Finish with one table: \emph{Layer $|$ Change $|$ S/D $|$ Why}.

\subsection{The layers of a processor (what you are allowed to change)}\label{sec:layers}
\small
\begin{longtable}{L{2.7cm}L{4.6cm}L{8.3cm}}
\toprule
\tblhead{\textbf{Layer} & \textbf{What it does} & \textbf{Concepts}}\\
\midrule
\endhead
Work distribution (software) & Splits the program and maps it onto PEs & \cp{W1}--\cp{W9}; compiler: \cp{I1} \cp{I2} \cp{I8} \cp{M13} \cp{M14} \\
Front end & Fetches the right instructions & \cp{P6} \cp{P7} \cp{P8} \cp{P9} \cp{P10} \cp{P11} \cp{P12} \cp{P13} \cp{M16} \\
Pipeline and execution & Executes instructions & \cp{P1} \cp{P2} \cp{P4} \cp{P5}; \cp{I3} \cp{I4} \cp{I5} \cp{I6} \cp{I7} \cp{I9} \\
Threads per core & Hides stalls using other threads & \cp{T1} \cp{T2} \cp{T3} \\
Data-parallel units & One instruction, many data & \cp{D1} \cp{D2} \cp{D3} \cp{D4} \cp{D5} \cp{D6} \cp{D7} \cp{D8} \\
Number/kind of PEs & More or different workers & \cp{T4} \cp{T5} \cp{T6} \cp{E3} \cp{W9} \\
Caches & Keep data close & \cp{M1}--\cp{M16}, \cp{M20}, \cp{T6} \\
Main memory, VM & Bandwidth, latency, protection & \cp{M17} \cp{M18} \cp{M19} \\
Interconnect, coherence, sync & Lets PEs share and order data & \cp{C1}--\cp{C9} \\
Power and reliability & Energy budget, uptime & \cp{E1} \cp{E2} \cp{E3} \cp{E4} \\
\bottomrule
\end{longtable}
\normalsize
